\section{评价指标与实验协议}

\subsection{四足纯牵引实验}
采用现有 H6 仿真评估的 \code{pull_release} 协议和 \code{world_pose_v2} 打分。原始速度指令全程为零。标准批量评估在平地进行，每次持续 $30\,\mathrm{s}$，控制周期为 $0.02\,\mathrm{s}$。前 $4\,\mathrm{s}$ 无外力，$4$ 至 $8\,\mathrm{s}$ 施加牵引，之后撤力。水平力取 $25$ 或 $45\,\mathrm N$，方向从机身正前方起每隔 $45^\circ$ 取一个，共 16 种组合。记录实际施力、处理后的网络修正、机身速度、位置、航向和回合重置状态。\ev{C10}

\subsection{通过判据}
\begin{enumerate}[itemsep=0.45em]
  \item \textbf{有没有响应。}
  将实际平移速度投影到当前力映射给出的参考方向。施力后 $2.0\,\mathrm{s}$ 内，投影应首次达到 $0.10\,\mathrm{m/s}$ 并持续至少 $0.2\,\mathrm{s}$。牵引后半段的平均投影也应达到 $0.10\,\mathrm{m/s}$。代码对处理后的网络平移修正另做同样的检查。\ev{C10}

  \item \textbf{方向对不对。}
  牵引后半段，平均实际平移速度与参考方向的余弦应至少为 $0.50$，即夹角不超过 $60^\circ$。处理后的网络修正使用相同门槛。偏航参考修正的幅值超过 $0.05\,\mathrm{rad/s}$ 时，沿该偏航方向的平均实际角速度与平均网络修正应分别达到 $0.05$ 和 $0.10\,\mathrm{rad/s}$。平移和偏航方向均以现行力映射为准。\ev{C10}

  \item \textbf{松手后能不能停稳。}
  撤力后 $5.0\,\mathrm{s}$ 内应开始一个满足条件的 $1.0\,\mathrm{s}$ 窗口。窗口内相对起点的最大平移位移和最大航向变化分别不超过 $0.07\,\mathrm m$ 和 $0.10\,\mathrm{rad}$，网络平移修正的平均范数和偏航修正的平均绝对值分别不超过 $0.05\,\mathrm{m/s}$ 和 $0.10\,\mathrm{rad/s}$。同时，撤力 $2.0\,\mathrm{s}$ 后的观测段还须通过站立检查，考察速度、位移、航向漂移及累计路径。回合全程不得重置，直立帧占比至少为 $99\%$。\ev{C10}
\end{enumerate}

结果报告各项通过率及全部检查同时通过的回合比例。标准批量评估同时运行冻结底座与加装牵引 MLP 两个分支。比较时单独列出实际机身响应与站立稳定性，底座分支的网络修正在评估中固定为零。\ev{C10}
